\section{Análisis de stakeholders del \texorpdfstring{\ac{GRID}}{GRID}}

Con el propósito de identificar los actores relacionados con la infraestructura tecnológica del Grupo de Investigación en Redes, Información y Distribución (\ac{GRID}), se llevó a cabo un análisis de stakeholders, cuyos resultados precisaron los diferentes roles, intereses y niveles de influencia que presentan los actores frente a la gestión y utilización de los servicios tecnológicos considerados en el proyecto.

Los principales stakeholders identificados corresponden al Grupo de Investigación \ac{GRID}, los docentes del programa de Ingeniería de Sistemas y Computación y los estudiantes del mismo programa. Estos actores mantienen una relación directa con la infraestructura, aunque con diferentes niveles de participación, influencia e interés.

El Grupo de Investigación \ac{GRID} constituye el principal beneficiario y actor con capacidad de decisión frente al proyecto, debido a que es responsable de la infraestructura tecnológica y de los servicios sobre los cuales se plantea la solución. Su participación resulta necesaria para establecer las necesidades de gestión de usuarios, definir las condiciones de acceso a los recursos y determinar las características que debe cumplir la solución implementada.

Los docentes de Ingeniería de Sistemas y Computación representan usuarios de la infraestructura tecnológica, principalmente en el desarrollo de actividades académicas e investigativas que requieren el acceso a los servicios proporcionados por el \ac{GRID}. Por su relación con estos recursos, presentan interés en disponer de mecanismos que permitan un acceso adecuado a los servicios según las autorizaciones correspondientes.

Por su parte, los estudiantes de Ingeniería de Sistemas y Computación corresponden a usuarios finales de los recursos tecnológicos disponibles. Su participación se relaciona con el uso de los servicios para actividades académicas, prácticas y de formación, por lo que constituyen un grupo directamente afectado por los mecanismos establecidos para gestionar las cuentas y controlar el acceso.

La \autoref{tab:stakeholders} presenta la caracterización de estos stakeholders, considerando su rol, relación con el proyecto, impacto, poder de influencia, interés y compromiso. Esta identificación permite establecer las responsabilidades y expectativas de los principales actores involucrados en el contexto de la solución propuesta.

\begin{table}[H]
	\centering
	\caption{Análisis de stakeholders}
	\label{tab:stakeholders}

	\renewcommand{\arraystretch}{1.25}
	\setlength{\tabcolsep}{4pt}
	\fontsize{7pt}{7pt}\selectfont

	\begin{tabularx}{\textwidth}{|
		>{\raggedright\arraybackslash}p{2.0cm}|
		>{\raggedright\arraybackslash}p{1.7cm}|
		>{\raggedright\arraybackslash}X|
		>{\raggedright\arraybackslash}p{1.5cm}|
		>{\raggedright\arraybackslash}X|
		>{\raggedright\arraybackslash}X|
		>{\raggedright\arraybackslash}p{2.0cm}|}

		\hline
		\textbf{Interesado}                                                                                                                                          &
		\textbf{Rol}                                                                                                                                                 &
		\textbf{Relación}                                                                                                                                            &
		\textbf{Impacto}                                                                                                                                             &
		\textbf{Poder de influencia}                                                                                                                                 &
		\textbf{Interés}                                                                                                                                             &
		\textbf{Compromiso}
		\\
		\hline

		Docentes de Ingeniería de Sistemas                                                                                                                           &
		Usuarios clave                                                                                                                                               &
		Utilizarán los servicios tecnológicos del \acs{GRID} en actividades académicas e investigativas que requieran acceso a los recursos de la infraestructura          &
		Medio-Alto                                                                                                                                                   &
		Medio, pueden influir en la adopción mediante el uso de los servicios, aunque no tienen responsabilidad directa sobre las decisiones de la infraestructura   &
		Medio-Alto, requieren mecanismos de acceso que faciliten el uso de los servicios tecnológicos de acuerdo con los permisos establecidos                       &
		Medio, condicionado por la utilidad y facilidad de uso de los mecanismos de autenticación y acceso
		\\
		\hline

		Estudiantes de Ingeniería de Sistemas                                                                                                                        &
		Usuarios finales                                                                                                                                             &
		Utilizarán los servicios tecnológicos disponibles para actividades académicas, prácticas y de formación                                                      &
		Bajo-Medio                                                                                                                                                   &
		Bajo, no participan directamente en las decisiones sobre la infraestructura, aunque su experiencia de uso permite valorar el funcionamiento de la solución   &
		Medio, debido a la necesidad de acceder a los recursos tecnológicos requeridos para sus actividades académicas                                               &
		Bajo-Medio, relacionado con la disponibilidad y facilidad de acceso a los servicios
		\\
		\hline

		Grupo de Investigación \acs{GRID}                                                                                                                                  &
		Beneficiario principal                                                                                                                                       &
		Proporciona y administra la infraestructura tecnológica, define las necesidades de gestión de usuarios y establece las condiciones de acceso a los servicios &
		Alto                                                                                                                                                         &
		Alto, tiene capacidad para definir las condiciones de implementación y tomar decisiones sobre la gestión de la infraestructura                               &
		Alto, busca mejorar la administración de las cuentas de usuario, centralizar la gestión de identidades y facilitar el control de acceso a sus servicios      &
		Alto, debido a su responsabilidad directa sobre la infraestructura y a los beneficios esperados de la solución
		\\
		\hline
	\end{tabularx}
\end{table}

